\paragraph{A dislocation above the inclusion.} A slab of the same cross-section,
$126.0\times124.0\times125.5$~\AA{} (117\,363 atoms), holds an Al$_{13}$Fe$_4$
sphere of radius 25~\AA{} in its lower half; the slip plane is 12~\AA{} above
the top of the sphere, and as built the sphere centre is 15~\AA{} ahead of the
line along $x$. The inclusion is held as in the field calculations, and the
blocks are rigid, as in the calculation without an inclusion. The matrix is relaxed by conjugate gradients to a force norm
of $10^{-3}$~eV/\AA{} without annealing, and the top block is then moved along
$z$ until the mean $\sigma_{zz}$ of the matrix is close to zero. During this
relaxation the line moves «DISL_RELAX»~\AA{} away from the inclusion and stops
«DISL_D»~\AA{} from the sphere centre along $x$. The line position is the mean
position of the two partial dislocations found by DXA; their separation is
«DISL_SPLIT»~\AA{} in every state. The inclusion atoms are then given the strain
$+\varepsilon$ and the matrix is relaxed again: the line moves a further
«DISL_SHIFT»~\AA{} away from the inclusion. After the inclusion atoms are put
back and the matrix is relaxed once more, the line is «DISL_REV»~\AA{} from
its first position.

Shearing these slabs did not give a separate measure of the effect of the
field. Even without load the top and bottom blocks carry shear stresses of
«BLOCK_TOP» and «BLOCK_BOT»~MPa; the held inclusion transmits the difference,
so the stress at which the line breaks away is set mainly by the forces that
hold the inclusion in place, and the contribution of the inclusion strain
cannot be separated in these runs.
